% ch3.tex -- Module 6 ESSENTIALS, Chapter 3: Electromagnetic induction (flux, Faraday, Lenz)
% Condensed from src/mod6/ch3.tex (final full book). Knowledge points only: no questions,
% no worked examples, no exam tips. Same notation and facts as the full chapter.
% Figure files: figures/ch3_falling.pdf (+ .py), copied from the full book. All other figures
% are TikZ copied (some shrunk) from the full chapter.
\chapter{Electromagnetic Induction: Flux, Faraday and Lenz}\label{ch3:chap}
\covers{M6.IQ3.01, M6.IQ3.02}
% Local TikZ helpers (chapter prefix): a small open loop with one arrowhead.
% \chThreeACW{x}{y}{r} = anticlockwise loop, \chThreeCW{x}{y}{r} = clockwise loop.
\newcommand\chThreeACW[3]{%
  \draw[line width=1.2pt, hscink] ($(#1,#2)+(110:#3)$) arc (110:400:#3);
  \draw[line width=1.2pt, hscink, -{Latex[length=2.4mm, width=2.0mm]}]
    ($(#1,#2)+(40:#3)+(310:0.05)$) -- ($(#1,#2)+(40:#3)+(130:0.19)$);}
\newcommand\chThreeCW[3]{%
  \draw[line width=1.2pt, hscink] ($(#1,#2)+(70:#3)$) arc (70:-220:#3);
  \draw[line width=1.2pt, hscink, -{Latex[length=2.4mm, width=2.0mm]}]
    ($(#1,#2)+(140:#3)+(230:0.05)$) -- ($(#1,#2)+(140:#3)+(50:0.19)$);}

\begin{summarybox}{In this chapter}
\begin{itemize}
  \item Magnetic flux, $\Phi = B_\parallel A = BA\cos\theta$, and how it can change.
  \item Faraday's law: a changing flux induces an emf.
  \item Lenz's law: the induced current opposes the change (conservation of energy).
  \item Induction in moving rods, solenoids, metal plates and two nearby coils.
\end{itemize}
\end{summarybox}

%% ===================================================================
\section{Magnetic flux}\label{ch3:sec:flux}
\covers{M6.IQ3.01}

\begin{bigidea}
Magnetic flux tells you how much magnetic field passes through a loop. It
depends on the field strength, the area and the angle.
\end{bigidea}

\begin{keywords}
  \kw{Magnetic flux ($\Phi$)}{how much magnetic field passes through an area.
    Unit: weber (Wb).}
  \kw{Normal}{an imaginary arrow at \ang{90} to the flat surface of a loop.}
\end{keywords}

\begin{figure}
  \centering
  \begin{tikzpicture}[x=1cm, y=1cm]
    % ---------- (a) theta = 0
    \begin{scope}
      \foreach \y in {0.2, 0.75, 1.8, 2.35} \draw[field] (0,\y) -- (3.6,\y);
      \node[diagramlabel, text=hscfield, right] at (3.6,2.35) {$B$};
      \coordinate (C) at (1.5,1.275);
      \draw[line width=2.4pt, hscbrand] ($(C)+(90:0.7)$) -- ($(C)+(270:0.7)$);
      \draw[->, line width=1.1pt, hscink] (C) -- ++(0:1.0)
        node[pos=1, right, font=\footnotesize, inner sep=2pt] {normal};
      \node[align=center, font=\small] at (1.8,-0.5)
        {(a) $\theta = \ang{0}$\\ $\Phi = BA$ (largest)};
    \end{scope}
    % ---------- (b) theta = 50 deg
    \begin{scope}[xshift=4.75cm]
      \foreach \y in {0.2, 0.75, 1.8, 2.35} \draw[field] (0,\y) -- (3.6,\y);
      \node[diagramlabel, text=hscfield, right] at (3.6,2.35) {$B$};
      \coordinate (C) at (1.5,1.275);
      \draw[guide] (C) -- ++(0:1.1);
      \draw[line width=2.4pt, hscbrand] ($(C)+(140:0.7)$) -- ($(C)+(320:0.7)$);
      \draw[->, line width=1.1pt, hscink] (C) -- ++(50:0.9)
        node[anchor=south west, font=\footnotesize, inner sep=1pt] {normal};
      \draw[line width=0.7pt, hscink] ($(C)+(0:0.45)$) arc (0:50:0.45);
      \node[font=\small, inner sep=0.5pt] at ($(C)+(22:0.68)$) {$\theta$};
      \draw[line width=0.8pt, hscmuted] ($(C)+(320:0.55)$) arc (320:360:0.55);
      \node[font=\footnotesize, text=hscmuted, inner sep=1pt, anchor=west]
        at ($(C)+(338:0.8)$) {$\ang{90} - \theta$};
      \node[align=center, font=\small] at (1.8,-0.5)
        {(b) tilted\\ $\Phi = BA\cos\theta$};
    \end{scope}
    % ---------- (c) theta = 90 deg
    \begin{scope}[xshift=9.5cm]
      \foreach \y in {0.2, 0.75, 1.8, 2.35} \draw[field] (0,\y) -- (3.6,\y);
      \node[diagramlabel, text=hscfield, right] at (3.6,2.35) {$B$};
      \coordinate (C) at (1.2,1.275);
      \draw[line width=2.4pt, hscbrand] ($(C)+(180:0.7)$) -- ($(C)+(0:0.7)$);
      \draw[->, line width=1.1pt, hscink] (C) -- ++(90:0.9)
        node[pos=0.89, anchor=east, font=\footnotesize, inner sep=1.5pt, fill=white] {normal};
      \node[align=center, font=\small] at (1.8,-0.5)
        {(c) $\theta = \ang{90}$: plane\\ parallel to $B$, \ $\Phi = 0$};
    \end{scope}
  \end{tikzpicture}
  \caption{A flat loop in a uniform field, seen from the side (thick navy
    line).}\label{ch3:fig:angle}
\end{figure}

\begin{itemize}
  \item The \term{plane} of a loop is the flat surface inside it, like the skin of a
        drum. $\theta$ is measured from the \emph{normal}, not from the plane
        (\cref{ch3:fig:angle}).
  \item Only the part of $B$ along the normal, $B_\parallel = B\cos\theta$, passes
        through. For an angle to the plane, use $\theta = \ang{90} - \text{that angle}$.
  \item $\qty{1}{Wb} = \qty{1}{T.m^2}$. With $B$ perpendicular to the plane,
        $B = \dfrac{\Phi}{A}$, the flux per square metre. So $B$ is also called the
        \term{magnetic flux density}.
  \item Reversing the field (or turning the loop over) takes the flux from $+BA$ to
        $-BA$. So $\lvert\Delta\Phi\rvert = 2BA$, not zero.
\end{itemize}


\begin{formulacard}{Magnetic flux}
  \cardformula{\Phi = B_\parallel A = BA\cos\theta}
  \begin{symbols}
    \sym{$\Phi$}{magnetic flux through the loop}{\si{Wb} ($= \si{T.m^2}$)}
    \sym{$B_\parallel$}{the part of $B$ along the normal of the loop}{\si{T}}
    \sym{$B$}{magnetic field strength (flux density)}{\si{T}}
    \sym{$A$}{area of the loop that is inside the field}{\si{m^2}}
    \sym{$\theta$}{angle between $B$ and the \emph{normal} to the loop}{\si{\degree}
      (SI unit: rad)}
  \end{symbols}
  \whentouse{a flat loop or coil in a uniform field. $\Phi$ is for \emph{one} turn.}
  \trap{plane \emph{parallel} to the field: $\theta = \ang{90}$, so $\Phi = 0$.
    Change \si{cm^2} to \si{m^2} ($\times\,10^{-4}$).}
\end{formulacard}

\subsection{Three ways to change the flux}
\begin{itemize}
  \item \textbf{Change $B$}: move a magnet closer or away; change the current in an
        electromagnet; switch the field on, off, or reverse it.
  \item \textbf{Change $A$}: slide the loop into or out of the field; squash it.
        \textbf{Change $\theta$}: rotate the loop, as in a generator.
  \item A loop sliding \emph{inside} a large uniform field keeps the same flux: no emf.
\end{itemize}

%% ===================================================================
\section{Faraday's law}\label{ch3:sec:faraday}
\covers{M6.IQ3.02}

\begin{bigidea}
Only a \emph{changing} flux induces an emf. A steady flux induces nothing, however
big it is.
\end{bigidea}

\begin{itemize}
  \item \term{Electromagnetic induction}: making (\emph{inducing}) an emf by changing flux.
  \item \term{emf} ($\varepsilon$, electromotive force): the energy given to each
        coulomb of charge (not a force). Unit: volt (V).
  \item A faster magnet, a stronger magnet or more turns give a bigger emf.
  \item \textbf{Faraday's law}: the induced emf is proportional to the rate of
        change of magnetic flux through the circuit.
  \item The $N$ turns add their emfs, like cells in a battery. Closed circuit: an
        \term{induced current} $I = \dfrac{\varepsilon}{R}$ flows. Open circuit: an
        emf, but no current.
  \item On a $\Phi$--$t$ graph, $\lvert\varepsilon\rvert = N \times \text{gradient}$.
        Steep: big emf. Flat: zero. Rising and falling parts give opposite signs.
\end{itemize}

\begin{formulacard}{Faraday's law of induction}
  \cardformula{\varepsilon = -N\frac{\Delta\Phi}{\Delta t}}
  \begin{symbols}
    \sym{$\varepsilon$}{average induced emf}{\si{V}}
    \sym{$N$}{number of turns in the coil}{no unit}
    \sym{$\Delta\Phi$}{change in flux through one turn:
      $\Phi_\mathrm{final} - \Phi_\mathrm{initial}$}{\si{Wb}}
    \sym{$\Delta t$}{time taken for the change}{\si{s}}
  \end{symbols}
  \whentouse{the size of the emf for any change of flux.}
  \trap{$\Delta\Phi$ is the \emph{change} in flux: a huge steady flux gives zero emf.
    $\si{Wb.s^{-1}} = \si{V}$.}
\end{formulacard}

%% ===================================================================
\section{Lenz's law: the direction of the current}\label{ch3:sec:lenz}
\covers{M6.IQ3.02}

\begin{bigidea}
An induced current always fights the change that made it. \term{Lenz's law}
(learn it): \emph{the induced current flows in the direction that makes a magnetic
field opposing the change in flux that produced it.}
\end{bigidea}

\begin{itemize}
  \item Lenz's law is the minus sign in Faraday's law. It follows from conservation
        of energy. Open circuit: the emf acts the way that would drive such a current.
  \item The induced field opposes the \emph{change}, not the field. Flux increasing: it
        points opposite to the original field. Flux decreasing: the \emph{same} way.
  \item \term{Right-hand grip rule}: fingers curl with the current, thumb gives the
        field inside. Anticlockwise (as you see it): an N face. Clockwise: an S face.
  \item Steps for any direction: (1) find which way the field points and whether the
        flux grows or shrinks; (2) Lenz's law gives the induced field; (3) the grip
        rule gives the current. Clockwise or anticlockwise has a meaning only when the
        viewing side is named (for example, seen from the magnet).
  \item \term{Pole method} (\cref{ch3:fig:magnetcoil}): the near face of the coil becomes
        the \emph{same} pole as an approaching pole (it repels), and the
        \emph{opposite} pole to a leaving pole (it attracts).
\end{itemize}

\begin{figure}
  \centering
  \begin{tikzpicture}[x=1cm, y=1cm]
    \foreach \xs in {0, 7.1} {
    \begin{scope}[xshift=\xs cm]
      \fill[hscvel!22] (0,0) rectangle (0.8,0.5);
      \fill[hscforce!22] (0.8,0) rectangle (1.6,0.5);
      \draw[line width=0.8pt, hscink] (0,0) rectangle (1.6,0.5);
      \node[font=\small\bfseries] at (0.4,0.25) {S};
      \node[font=\small\bfseries] at (1.2,0.25) {N};
      \foreach \x in {2.6, 2.85, 3.1, 3.35, 3.6}
        \draw[line width=1pt, hscink] (\x,0.25) ellipse (0.1 and 0.6);
    \end{scope}}
    % ----- (a) approaching
    \draw[velocity] (0.3,0.85) -- (1.3,0.85) node[right, text=hscvel] {$v$};
    \draw[->, hscfield, line width=1.2pt] (3.95,0.25) -- (2.05,0.25);
    \node[font=\footnotesize, text=hscfield, anchor=west] at (3.95,0.25) {induced $B$};
    \node[font=\footnotesize\bfseries, text=hscforce] at (3.1,1.15) {near face: N};
    \node[font=\footnotesize, align=center] at (2.2,-0.95)
      {(a) N pole approaching: the near\\ face becomes N and repels.};
    \chThreeACW{1.2}{-1.9}{0.36}
    \node[font=\footnotesize, align=left, anchor=west] at (1.75,-1.9)
      {seen from the magnet:\\ anticlockwise};
    % ----- (b) moving away
    \begin{scope}[xshift=7.1cm]
      \draw[velocity] (1.3,0.85) -- (0.3,0.85) node[left, text=hscvel] {$v$};
      \draw[->, hscfield, line width=1.2pt] (2.05,0.25) -- (3.95,0.25);
      \node[font=\footnotesize, text=hscfield, anchor=west] at (3.95,0.25) {induced $B$};
      \node[font=\footnotesize\bfseries, text=hscvel] at (3.1,1.15) {near face: S};
      \node[font=\footnotesize, align=center] at (2.2,-0.95)
        {(b) N pole moving away: the near\\ face becomes S and attracts.};
      \chThreeCW{1.2}{-1.9}{0.36}
      \node[font=\footnotesize, align=left, anchor=west] at (1.75,-1.9)
        {seen from the magnet:\\ clockwise};
    \end{scope}
  \end{tikzpicture}
  \caption{The pole method. The coil always opposes the motion of the
    magnet.}\label{ch3:fig:magnetcoil}
\end{figure}

\subsection{Lenz's law and conservation of energy}
\begin{itemize}
  \item If a coil \emph{attracted} an approaching magnet, the magnet would speed up
        and gain energy from nothing. That breaks the \term{law of conservation of
        energy}: energy cannot be created or destroyed.
  \item So the coil pushes back: work done against it $\to$ electrical energy $\to$
        heat. Open circuit: no current, no push back.
\end{itemize}

%% ===================================================================
\section{Motional emf: a straight wire moving in a field}\label{ch3:sec:motional}
\covers{M6.IQ3.02}

\begin{bigidea}
A metal rod moving across a magnetic field has an emf between its ends.
\end{bigidea}

\begin{sidebyside}[0.3]
  \centering
  \begin{tikzpicture}[x=0.8cm, y=0.8cm]
    \foreach \x in {2.1, 2.7} \foreach \y in {0.3, 0.9, 2.1, 2.7}
      \node[text=hscfield, font=\small] at (\x,\y) {$\times$};
    \foreach \y in {0.3, 0.9, 1.5, 2.1, 2.7}
      \node[text=hscfield, font=\small] at (3.3,\y) {$\times$};
    \foreach \x in {0.3, 0.9} \foreach \y in {0.3, 0.9, 2.7}
      \node[text=hscfield, font=\small] at (\x,\y) {$\times$};
    \draw[line width=3.2pt, hscmuted] (1.5,0.15) -- (1.5,2.85);
    \node[draw, circle, inner sep=0.6pt, font=\footnotesize\bfseries, fill=white]
      at (1.5,3.15) {$+$};
    \node[draw, circle, inner sep=0.6pt, font=\footnotesize\bfseries, fill=white]
      at (1.5,-0.15) {$-$};
    \draw[velocity] (1.62,1.2) -- (2.5,1.2);
    \node[text=hscvel, font=\small, fill=white, inner sep=1pt, anchor=west]
      at (2.52,1.2) {$v$};
    \draw[force] (1.2,1.25) -- (1.2,2.35);
    \node[font=\footnotesize, text=hscforce, fill=white, inner sep=1pt, anchor=east]
      at (1.12,1.8) {$F$ on $+q$};
  \end{tikzpicture}
  \captionof{figure}{A rod alone, $B$ into the page.}\label{ch3:fig:rod}
\nextside
  \begin{itemize}
    \item \term{Right-hand palm rule}: fingers along $B$, thumb along $v$ of a
          \emph{positive} charge, palm pushes the way of the force.
    \item Rod moving right (\cref{ch3:fig:rod}): positive charge is pushed up, free
          electrons (they can move in the metal) down. The top end becomes $+$, like a
          battery: \term{motional emf}.
    \item The build-up stops when the electric and magnetic forces balance.
  \end{itemize}
\end{sidebyside}

\begin{sidebyside}[0.43]
  \centering
  \begin{tikzpicture}[x=0.75cm, y=0.75cm]
    \fill[hscvel!12] (3.0,0) rectangle (3.8,3.0);
    \foreach \x in {1.65, 2.4, 4.6, 5.4} \foreach \y in {0.5, 1.5, 2.5}
      \node[text=hscfield, font=\small] at (\x,\y) {$\times$};
    \foreach \y in {0.5, 2.5} \node[text=hscfield, font=\small] at (0.9,\y) {$\times$};
    \foreach \y in {0.5, 2.5} \node[text=hscfield, font=\small] at (3.4,\y) {$\times$};
    \draw[line width=1.5pt, hscink] (0,0) -- (6.0,0);
    \draw[line width=1.5pt, hscink] (0,3.0) -- (6.0,3.0);
    \draw[line width=1.2pt, hscink] (0,0) -- (0,1.0);
    \draw[line width=1.2pt, hscink] (0,2.0) -- (0,3.0);
    \draw[line width=1.2pt, hscink, fill=white] (-0.16,1.0) rectangle (0.16,2.0);
    \node[font=\small, anchor=west] at (0.2,1.5) {$R$};
    \draw[line width=1pt, hscmuted, dashed] (3.0,-0.2) -- (3.0,3.2);
    \draw[line width=3.2pt, hscmuted] (3.8,-0.2) -- (3.8,3.2);
    \draw[velocity] (3.92,2.0) -- (4.9,2.0);
    \node[text=hscvel, font=\small, fill=white, inner sep=1pt, anchor=west]
      at (4.92,2.0) {$v$};
    \draw[force] (3.68,1.0) -- (2.7,1.0);
    \node[font=\footnotesize, text=hscforce, fill=white, inner sep=1pt, anchor=east]
      at (2.65,1.0) {$F = lIB$};
    \draw[->, line width=1pt, hscink] (4.05,0.2) -- (4.05,0.95)
      node[right, font=\small] {$I$};
    \draw[->, line width=1pt, hscink] (2.2,3.28) -- (1.2,3.28) node[left, font=\small] {$I$};
    \draw[->, line width=1pt, hscink] (1.2,-0.28) -- (2.2,-0.28) node[right, font=\small] {$I$};
    \draw[->, line width=1pt, hscink] (-0.45,2.0) -- (-0.45,1.0)
      node[midway, left, font=\small] {$I$};
    \draw[<->, line width=0.6pt, hscmuted] (6.35,0) -- (6.35,3.0)
      node[midway, right, text=hscink, font=\small] {$l$};
    \draw[<->, line width=0.6pt, hscmuted] (3.0,-0.62) -- (3.8,-0.62);
    \node[font=\small, below] at (3.4,-0.64) {$v\Delta t$};
  \end{tikzpicture}
  \captionof{figure}{The rod on rails. It sweeps the shaded area in time
    $\Delta t$.}\label{ch3:fig:rails}
\nextside
  \begin{itemize}
    \item On rails (\cref{ch3:fig:rails}) the current flows up the rod: anticlockwise.
    \item The force $F = lIB$ on the rod \emph{opposes} the motion (Lenz). The push
          does work, which becomes heat in $R$.
    \item Faraday's view: the rod sweeps out area, so the flux grows. Both views agree.
  \end{itemize}
\end{sidebyside}

\begin{itemize}
  \item Rod moving along the field lines, or along its own length: $\varepsilon = 0$.
  \item Induction: motion causes a current. Motor effect: a current causes motion.
\end{itemize}

\begin{formulacard}[notonsheet]{Motional emf}
  \cardformula{\varepsilon = lvB}
  \begin{symbols}
    \sym{$\varepsilon$}{emf between the ends of the rod}{\si{V}}
    \sym{$l$}{length of the rod inside the field}{\si{m}}
    \sym{$v$}{speed of the rod across the field}{\si{m.s^{-1}}}
    \sym{$B$}{magnetic field strength}{\si{T}}
  \end{symbols}
  \whentouse{rod, velocity and field all at \ang{90} to each other.}
  \derive{in time $\Delta t$ the rod sweeps an area $\Delta A = l\,v\Delta t$
    (\cref{ch3:fig:rails}), so $\Delta\Phi = B\,l\,v\Delta t$. Faraday's law with
    $N = 1$: $\lvert\varepsilon\rvert = \dfrac{\Delta\Phi}{\Delta t} = lvB$.}
  \trap{$l$ is the length \emph{inside the field}, not the whole wire.}
\end{formulacard}

%% ===================================================================
\section{Solenoids, metal plates and two coils}\label{ch3:sec:solenoids}
\covers{M6.IQ3.02}

Any change in flux induces an emf, whatever causes it. Magnets in a
\term{solenoid} (a long coil), metal plates and nearby coils follow the same laws.

\subsection{A magnet moving through a solenoid}
\begin{itemize}
  \item Approaching: the flux increases, so an emf is induced one way. Centre: the flux
        is largest and not changing, so $\varepsilon = 0$. Leaving: the flux decreases,
        so the emf is the other way (\cref{ch3:fig:falling}).
\end{itemize}

\begin{sidebyside}[0.64]
  \centering
  \includegraphics[width=\linewidth]{ch3_falling.pdf}
  \captionof{figure}{Model: a magnet falling through a coil.}\label{ch3:fig:falling}
\nextside
  \begin{itemize}[leftmargin=1.1em]
    \item A faster magnet gives taller, narrower pulses. A falling magnet leaves
          faster, so pulse 2 is taller and narrower.
    \item Each pulse has the same area, $N\lvert\Delta\Phi\rvert$.
  \end{itemize}
\end{sidebyside}

\begin{itemize}
  \item Closed circuit: the coil repels the magnet above it and attracts it below.
        Some gravitational potential energy becomes heat, so the magnet falls slower.
\end{itemize}

\subsection{Metal plates: eddy currents}
\begin{itemize}
  \item \term{Eddy currents}: loops of current induced inside solid metal, only where
        the flux through it changes. Lenz's law gives their direction.
  \item Plate moving right, field into the page: anticlockwise where the plate enters
        the field; none when fully inside; clockwise where it leaves.
  \item The force on each eddy loop opposes the motion. The plate slows, and its
        kinetic energy becomes heat: \term{magnetic braking}. Uses: brakes, induction
        cooktops, metal detectors. Problem: energy loss in transformer cores.
\end{itemize}

\subsection{One solenoid near another}
\begin{itemize}
  \item \term{Mutual induction}: a changing current in the \term{primary coil} (joined
        to the supply) induces an emf in a nearby \term{secondary coil}. A coil with a
        steady current, moved in or out, acts like a bar magnet.
  \item Switch closed: a brief emf; the secondary's induced field opposes the growing
        field. Kept closed: no emf. Opened: a brief emf the other way.
  \item A \term{soft iron} core (easily magnetised) guides more flux into the
        secondary: more emf.
  \item An \term{AC} supply (current keeps changing size and direction) gives an emf
        all the time: a transformer. A steady \term{DC} supply gives an emf only when
        switched.
  \item The secondary's energy comes from the supply (or your work, if a coil moves).
  \item \textbf{A coil wound on a long air-core solenoid} (all the primary's flux passes
        through it): each turn has $\Phi = BA$, with
        $B = \frac{\mu_0 N_\mathrm{p} I_\mathrm{p}}{L}$. Faraday's law then gives
        $\lvert\varepsilon_\mathrm{s}\rvert = N_\mathrm{s}\,\frac{\mu_0 N_\mathrm{p}
        A}{L}\,\frac{\lvert\Delta I_\mathrm{p}\rvert}{\Delta t}$ (not printed in this
        form). A faster change of $I_\mathrm{p}$ gives a bigger emf. $N_\mathrm{p}$,
        $N_\mathrm{s}$ = turns of the primary, of the secondary; $A$, $L$ = cross-section
        area and length of the primary.
\end{itemize}

%% ===================================================================
\section{Key experiments}\label{ch3:sec:experiments}
\covers{M6.IQ3.02}

\begin{experiment}{Faraday's induction experiments (1831)}
\aim{To show that a changing flux induces an emf.}
\method{A magnet was pushed into and pulled out of a coil joined to a galvanometer (a
  meter for very small currents). A wire was moved across a magnet's field. On two coils sharing an iron ring, a
  switch was closed and opened.}
\results{The needle flicked only while something changed. Opposite changes gave
  opposite flicks. Faster changes gave bigger flicks.}
\conclusion{An emf is induced only while the flux changes. Only relative motion
  matters.}
\significance{With Ørsted's work (1820), it linked electricity and magnetism.}
\end{experiment}

\begin{experiment}{Eddy currents: falling magnets and swinging plates}
\aim{To show that eddy currents oppose the motion that causes them.}
\method{A magnet fell through PVC, copper and slotted copper tubes. Solid and
  slotted plates swung between magnet poles.}
\results{Fall times: PVC \qty{0.41}{s} (free fall), slotted copper \qty{0.56}{s},
  copper \qty{3.57}{s}. The solid plate stopped fast; the slotted one swung on.}
\conclusion{Eddy currents in the copper oppose the fall (Lenz); the lost potential
  energy becomes heat. Slots break the current loops, so braking is weaker.}
\significance{Evidence for Lenz's law and energy conservation; used in magnetic
  brakes.}
\end{experiment}

%% ===================================================================
\begin{summarybox}{Remember}
\begin{itemize}
  \item $\Phi = BA\cos\theta$, $\theta$ from the normal. It changes if $B$, $A$ or
        $\theta$ changes.
  \item Only a \emph{changing} flux induces an emf:
        $\varepsilon = -N\dfrac{\Delta\Phi}{\Delta t}$.
  \item Lenz: the induced current opposes the \emph{change} in flux (energy is conserved).
  \item Moving rod: $\varepsilon = lvB$ (not on the sheet).
  \item Magnet through a coil: two opposite pulses. Two coils: an emf only while the
        primary current changes.
\end{itemize}
\end{summarybox}
